%! Author = lili
%! Date = 9/7/26

\section{Untersuchung der Genexpression mittels RT-PCR}

\subsection{Einführung}
Bestimmung der Konzentration und Reinheit von Nukleinsäuren



%\subsection{Isolierung von Gesamt-RNA}

\subsection{Photometrische Messung der RNA-Konzentration}

\ausschnitt{Die Konzentration einer DNA- oder RNA-Lösung wird üblicherweise
photometrisch bei einer Wellenlänge von 260 nm in einer Quarzkürvette
bestimmt. – Oder am Nanodrop \\
Aus der Extinktion der Lösung kann die Konzentration anhand des
folgenden Verhältnisses berechnet werden:
    \begin{itemize}
        \item $\mathrm{OD}_{260} = 1$ entspricht $50\,\mu\mathrm{g/mL}$ doppelsträngiger (ds) DNA
        \item $\mathrm{OD}_{260} = 1$ entspricht $40\,\mu\mathrm{g/mL}$ RNA
        \item $\mathrm{OD}_{260} = 1$ entspricht $33\,\mu\mathrm{g/mL}$ einzelsträngiger DNA
        \item $\mathrm{OD}_{260} = 1$ entspricht $20\,\mu\mathrm{g/mL}$ bei Oligonukleotiden
    \end{itemize}
Dieses gilt für Nukleinsäuren bei pH 7,0; die Messung sollte daher in einem
Puffer mit geringer Salzkonzentration (z.B. ein Tris)-Puffer mit pH 7,0)
    durchgeführt werden.\\
    Das Verhältnis der OD260 zur OD280 zeigt an, wie stark eine DNA-Lösung
    noch durch Alkohol und Reste von Proteinen verunreinigt ist.
}{\cite{folien32}}

\begin{prinz}
    Ein Verhältnis von 1,8 spricht für eine reine DNA-Isolierung \\
    Ein Verhältnis von 2,0 für eine reine RNA-Isolierung \\
    Ist die Nukleinsäure-Lösung mit Proteinen oder Phenol kontaminiert, so ist
    der Wert signifikant kleiner \cite{folien32}
\end{prinz}
\ausschnitt{Oft wird der Quotient OD230 zur OD260 bestimmt: Im Idealfall bei 0,45. Er ist
ein Indikator für weitere Verunreinigungen z.B. mit Polysacchariden ist.}{\cite{folien32}}



\kfrage{reinheit}
\kfrage{rtpcr}
\kfrage{rtpcrprimer}
\frage{rnaalspcr}
\frage{aktinkontrolle}
\frage{od260}
\kfrage{rnahäufigkeit}
\kfrage{genexpressiontranskript}
\kfrage{genexpressionmethoden}
\kfrage{genexpressionsanalyse_frage}
\kfrage{transkriptionbistranslation}
\kfrage{zwischentranskriptionundtranslation}
\kfrage{versuch_gewebeexpression}